%% TeXora template: Media Handout. Not slides: a tagged handout (an article) with a video, a GIF, a
%% picture and an animation, to test TeXora's packages in other editors. It compiles in any LaTeX
%% editor: TeXora, Overleaf, TeXstudio, TeXShop, VS Code with LaTeX Workshop.
%%
%% What it needs:
%%   - texoravideo.sty, texoraimage.sty and accessgate.sty BESIDE this file. TeXora puts them there
%%     when it creates this template or compiles it; for Overleaf, upload them with the .tex file.
%%   - media/sample-video.mp4, media/sample-loop.gif, media/sample-video-transcript.txt and
%%     media/test-card.png, which come with this template. Keep the media folder with the .tex file.
%%   - LaTeX from 2023 or later; LuaLaTeX gives the best tags.
%%
%% What to expect in another editor:
%%   - Before TeXora has compiled it once, a video shows as a box with its name (a LaTeX warning
%%     "[placeholder]", never an error). After that, its poster from texora-video/ shows. Keep
%%     texora-video/ and texora-images/ with the file when you copy or upload it.
%%   - Videos play in TeXora; other viewers show the poster. The animation plays in TeXora and Adobe
%%     Acrobat. The alt text, captions and transcripts are in every PDF.
%% Change the colours on the \definecolor lines below.
% !TEX program = lualatex
\IfFormatAtLeastTF{2025-06-01}
  {\DocumentMetadata{lang=en-US,pdfstandard=ua-2,tagging=on}}
  {\DocumentMetadata{lang=en-US,pdfstandard=ua-1,testphase={phase-III}}}
\documentclass[11pt]{article}
\usepackage[T1]{fontenc}
\usepackage{lmodern}
\usepackage[english]{babel}
\usepackage[margin=2.5cm]{geometry}
\usepackage{xcolor}
\usepackage{tikz}
\usepackage{animate}
\usepackage{hyperref}
\usepackage{texoravideo}
\usepackage{texoraimage}
\definecolor{forest}{RGB}{27,94,32}
\definecolor{moss}{RGB}{232,243,233}
\hypersetup{colorlinks,linkcolor=forest,urlcolor=forest,pdftitle={Media Handout},pdfauthor={Your Name}}

\title{Media Handout}
\author{Your Name\\ Your Institution}
\date{\today}

\begin{document}
\maketitle

\begin{abstract}
  A handout with moving pictures that stays usable with a screen reader. Each section shows one
  kind of media and the keys that describe it. Compile it in any editor to test the packages.
\end{abstract}

\section{A video}
The video below plays when clicked in TeXora. Its \texttt{alt} says what clicking does, its
\texttt{caption} what it shows, and its transcript, in a file of its own, everything in it.

\begin{center}
  \includevideo[width=0.6\textwidth][posterat=30,alt={Play video: a light beam turning, 6 seconds, no sound},caption={A light beam turns once around the centre.},transcriptfile=media/sample-video-transcript.txt]{media/sample-video.mp4}
\end{center}

\section{A GIF}
A GIF is written the same way. This one loops, so it does not start by itself: a click starts
and stops it.

\begin{center}
  \includevideo[height=4cm][loop,alt={Play animation: dots drifting from left to right, it loops},transcript={Small dots drift slowly from the left edge to the right edge, then start again. There is no text and no sound.}]{media/sample-loop.gif}
\end{center}

\section{A picture in any format}
\texttt{\textbackslash includeimage} takes PNG and JPEG, and also AVIF, WebP, HEIC and JPEG~XL,
which TeXora converts for LaTeX. It takes \texttt{alt} like \texttt{\textbackslash includegraphics}.

\begin{center}
  \includeimage[width=0.35\textwidth,alt={A television test card: eight vertical colour bars, a rainbow strip and a frame counter}]{media/test-card.png}
\end{center}

\section{An animation}
The animate package has no \texttt{alt} key, so \texttt{\textbackslash AccFigure} gives the
description. The \texttt{controls} show play and pause buttons.

\begin{center}
  \AccFigure{Five bars that grow one after another, from short on the left to tall on the right}{%
  \begin{animateinline}[controls,loop,alttext={Five bars growing one after another}]{2}
    \multiframe{5}{iN=1+1}{%
      \begin{tikzpicture}[y=0.5cm]
        \useasboundingbox (-0.3,-0.3) rectangle (5.4,5.3);
        \draw (-0.2,0) -- (5.2,0);
        \foreach \h [count=\i] in {1,2,3,4,5} {
          \ifnum\i>\iN\relax\else\fill[forest] (\i-0.8,0) rectangle (\i-0.2,\h);\fi
        }
      \end{tikzpicture}}
  \end{animateinline}}
\end{center}

\section{Transcripts}
\AccTranscripts

\end{document}
